% GENERATED FILE. Edit canonical lessons, then run npm run generate:books.
% canonical-source-hash: fnv1a64:d38dd06a9794a527
% canonical-lessons: KA-C416-muccalike, KA-C416-pratijne, KA-C416-pratiphala, KA-C416-koru, KA-C416-pramanavacana

\chapter{Promises and Rewards}
\label{ch:ka-promises-and-rewards}
\hlchaptermodality{416}

\begin{chapteropening}
\textbf{By the end of this chapter:} \emph{I can say a written undertaking, a vow, a reward, to ask formally and an oath of office.}

Complete the last lesson of chapter 416: I can say a written undertaking, a vow, a reward, to ask formally and an oath of office.
\end{chapteropening}

\section[\texorpdfstring{\kn{ಮುಚ್ಚಳಿಕೆ} (muccaḷike)}{ಮುಚ್ಚಳಿಕೆ (muccaḷike)}]{\kn{ಮುಚ್ಚಳಿಕೆ} (muccaḷike) — a written undertaking}

\label{lesson:KA-C416-muccalike}



\begin{quote}
Before the new one: say the Kannada for to sanction, then the Kannada for an official order.
\end{quote}

\subsection*{You'll want to know: \kn{ಮುಚ್ಚಳಿಕೆ}}
\textbf{\kn{ಮುಚ್ಚಳಿಕೆ}} — \emph{muccaḷike} — \textquotedblleft{}a written undertaking\textquotedblright{}.

\begin{quote}
\kn{ಮತ್ತೆ} \kn{ತಪ್ಪು} \kn{ಮಾಡುವುದಿಲ್ಲ} \kn{ಅಂತ} \kn{ಅವನು} \kn{ಮುಚ್ಚಳಿಕೆ} \kn{ಬರೆದುಕೊಟ್ಟನು}.
\end{quote}

\begin{quote}
\emph{matte tappu māḍuvudilla anta avanu muccaḷike baredukoṭṭanu.}
\end{quote}

\begin{quote}
He gave a written undertaking not to do wrong again.
\end{quote}

\kn{ಮುಚ್ಚಳಿಕೆ} \kn{ಬರೆದುಕೊಡು}, to sign an undertaking.

\begin{cousinweb}
\textbf{\kn{ಮುಚ್ಚಳಿಕೆ}} (\emph{muccaḷike}) is a Kannada word of the offices and courts whose origin is not settled: neither DEDR nor CDIAL lists it.
\end{cousinweb}

\begin{grammarlens}[title={What you've built}]
One more word for saying what happens if.
\end{grammarlens}

\subsection*{Your turn}
\begin{itemize}
\raggedright
  \item \emph{Say it:} \emph{muccaḷike}
  \item \emph{Say it:} \emph{muccaḷike}, once more
  \item \emph{Say it:} \emph{matte tappu māḍuvudilla anta avanu muccaḷike baredukoṭṭanu.} — He gave a written undertaking not to do wrong again.
  \item \emph{Recall:} say the Kannada for to sanction, then the Kannada for an official order, then say \emph{muccaḷike} again
\end{itemize}

\begin{tcolorbox}[breakable,colback=teal!4,colframe=teal!35!black,title={Before you move on}]
What does \kn{ಮುಚ್ಚಳಿಕೆ} mean? (\textquotedblleft{}A written undertaking\textquotedblright{}.) Say it once more.
\end{tcolorbox}
\section[\texorpdfstring{\kn{ಪ್ರತಿಜ್ಞೆ} (pratijñe)}{ಪ್ರತಿಜ್ಞೆ (pratijñe)}]{\kn{ಪ್ರತಿಜ್ಞೆ} (pratijñe) — a vow; a pledge}

\label{lesson:KA-C416-pratijne}



\begin{quote}
Before the new one: say the Kannada for an official order, then the Kannada for a written undertaking.
\end{quote}

\subsection*{You'll want to know: \kn{ಪ್ರತಿಜ್ಞೆ}}
\textbf{\kn{ಪ್ರತಿಜ್ಞೆ}} — \emph{pratijñe} — \textquotedblleft{}a vow; a pledge\textquotedblright{}.

\begin{quote}
\kn{ಇನ್ನು} \kn{ಮುಂದೆ} \kn{ಸಿಗರೇಟು} \kn{ಸೇದುವುದಿಲ್ಲ} \kn{ಅಂತ} \kn{ಅಪ್ಪ} \kn{ಪ್ರತಿಜ್ಞೆ} \kn{ಮಾಡಿದರು}.
\end{quote}

\begin{quote}
\emph{innu munde sigarēṭu sēduvudilla anta appa pratijñe māḍidaru.}
\end{quote}

\begin{quote}
Father vowed never to smoke again.
\end{quote}

\kn{ಪ್ರತಿಜ್ಞೆ} \kn{ಮಾಡು}, to vow.

\begin{cousinweb}
\textbf{\kn{ಪ್ರತಿಜ್ಞೆ}} (\emph{pratijñe}) is borrowed from Sanskrit \emph{pratijñā}, \textquotedblleft{}admission, promise, vow\textquotedblright{} (CDIAL 8563).
\end{cousinweb}

\begin{grammarlens}[title={What you've built}]
One more word for saying what happens if.
\end{grammarlens}

\subsection*{Your turn}
\begin{itemize}
\raggedright
  \item \emph{Say it:} \emph{pratijñe}
  \item \emph{Say it:} \emph{pratijñe}, once more
  \item \emph{Say it:} \emph{innu munde sigarēṭu sēduvudilla anta appa pratijñe māḍidaru.} — Father vowed never to smoke again.
  \item \emph{Recall:} say the Kannada for an official order, then the Kannada for a written undertaking, then say \emph{pratijñe} again
\end{itemize}

\begin{tcolorbox}[breakable,colback=teal!4,colframe=teal!35!black,title={Before you move on}]
What does \kn{ಪ್ರತಿಜ್ಞೆ} mean? (\textquotedblleft{}A vow; a pledge\textquotedblright{}.) Say it once more.
\end{tcolorbox}
\section[\texorpdfstring{\kn{ಪ್ರತಿಫಲ} (pratiphala)}{ಪ್ರತಿಫಲ (pratiphala)}]{\kn{ಪ್ರತಿಫಲ} (pratiphala) — a reward; a return}

\label{lesson:KA-C416-pratiphala}



\begin{quote}
Before the new one: say the Kannada for a written undertaking, then the Kannada for a vow.
\end{quote}

\subsection*{You'll want to know: \kn{ಪ್ರತಿಫಲ}}
\textbf{\kn{ಪ್ರತಿಫಲ}} — \emph{pratiphala} — \textquotedblleft{}a reward; a return\textquotedblright{}.

\begin{quote}
\kn{ಅವಳ} \kn{ಶ್ರಮಕ್ಕೆ} \kn{ಒಳ್ಳೆಯ} \kn{ಪ್ರತಿಫಲ} \kn{ಸಿಕ್ಕಿತು}.
\end{quote}

\begin{quote}
\emph{avaḷa śramakke oḷḷeya pratiphala sikkitu.}
\end{quote}

\begin{quote}
Her hard work brought a good reward.
\end{quote}

What comes back for what one gave.

\begin{cousinweb}
\textbf{\kn{ಪ್ರತಿಫಲ}} (\emph{pratiphala}) is borrowed from Sanskrit \emph{pratiphala}, \textquotedblleft{}a return, a requital\textquotedblright{} (Monier-Williams).
\end{cousinweb}

\begin{grammarlens}[title={What you've built}]
One more word for saying what happens if.
\end{grammarlens}

\subsection*{Your turn}
\begin{itemize}
\raggedright
  \item \emph{Say it:} \emph{pratiphala}
  \item \emph{Say it:} \emph{pratiphala}, once more
  \item \emph{Say it:} \emph{avaḷa śramakke oḷḷeya pratiphala sikkitu.} — Her hard work brought a good reward.
  \item \emph{Recall:} say the Kannada for a written undertaking, then the Kannada for a vow, then say \emph{pratiphala} again
\end{itemize}

\begin{tcolorbox}[breakable,colback=teal!4,colframe=teal!35!black,title={Before you move on}]
What does \kn{ಪ್ರತಿಫಲ} mean? (\textquotedblleft{}A reward; a return\textquotedblright{}.) Say it once more.
\end{tcolorbox}
\section[\texorpdfstring{\kn{ಕೋರು} (kōru)}{ಕೋರು (kōru)}]{\kn{ಕೋರು} (kōru) — to ask formally}

\label{lesson:KA-C416-koru}



\begin{quote}
Before the new one: say the Kannada for a vow, then the Kannada for a reward.
\end{quote}

\subsection*{You'll want to know: \kn{ಕೋರು}}
\textbf{\kn{ಕೋರು}} — \emph{kōru} — \textquotedblleft{}to ask formally\textquotedblright{}.

\begin{quote}
\kn{ನಾನು} \kn{ಕಚೇರಿಯಿಂದ} \kn{ಒಂದು} \kn{ವಾರದ} \kn{ರಜೆ} \kn{ಕೋರಿದೆ}.
\end{quote}

\begin{quote}
\emph{nānu kacēriyinda ondu vārada raje kōride.}
\end{quote}

\begin{quote}
I asked the office for a week of leave.
\end{quote}

Past (I …): \kn{ಕೋರಿದೆ}, \emph{kōride}.

Formal. \kn{ಕೋರಿಕೆ}, a request.

\begin{cousinweb}
\textbf{\kn{ಕೋರು}} (\emph{kōru}) is native Dravidian: DEDR 2232 lists Kannada \emph{kōṟu}, \emph{kōru}, \textquotedblleft{}to purpose, to wish, to desire\textquotedblright{}, beside Tamil \emph{kōru}, \textquotedblleft{}to request\textquotedblright{}.
\end{cousinweb}

\begin{grammarlens}[title={What you've built}]
One more word for saying what happens if.
\end{grammarlens}

\subsection*{Your turn}
\begin{itemize}
\raggedright
  \item \emph{Say it:} \emph{kōru}
  \item \emph{Say it:} \emph{kōru}, once more
  \item \emph{Say it:} \emph{nānu kacēriyinda ondu vārada raje kōride.} — I asked the office for a week of leave.
  \item \emph{Recall:} say the Kannada for a vow, then the Kannada for a reward, then say \emph{kōru} again
\end{itemize}

\begin{tcolorbox}[breakable,colback=teal!4,colframe=teal!35!black,title={Before you move on}]
What does \kn{ಕೋರು} mean? (\textquotedblleft{}To ask formally\textquotedblright{}.) Say it once more.
\end{tcolorbox}
\section[\texorpdfstring{\kn{ಪ್ರಮಾಣವಚನ} (pramāṇavacana)}{ಪ್ರಮಾಣವಚನ (pramāṇavacana)}]{\kn{ಪ್ರಮಾಣವಚನ} (pramāṇavacana) — an oath of office}

\label{lesson:KA-C416-pramanavacana}



\begin{quote}
Before the new one: say the Kannada for a reward, then the Kannada for to ask formally.
\end{quote}

\subsection*{You'll want to know: \kn{ಪ್ರಮಾಣವಚನ}}
\textbf{\kn{ಪ್ರಮಾಣವಚನ}} — \emph{pramāṇavacana} — \textquotedblleft{}an oath of office\textquotedblright{}.

\begin{quote}
\kn{ಹೊಸ} \kn{ಮಂತ್ರಿಗಳು} \kn{ಪ್ರಮಾಣವಚನ} \kn{ತೆಗೆದುಕೊಂಡರು}.
\end{quote}

\begin{quote}
\emph{hosa mantrigaḷu pramāṇavacana tegedukoṇḍaru.}
\end{quote}

\begin{quote}
The new ministers took the oath of office.
\end{quote}

\kn{ಪ್ರಮಾಣ} (proof, oath) + \kn{ವಚನ} (word).

\begin{cousinweb}
\textbf{\kn{ಪ್ರಮಾಣವಚನ}} (\emph{pramāṇavacana}) is built from Sanskrit \emph{pramāṇa}, \textquotedblleft{}a standard, a proof\textquotedblright{} (CDIAL 8723), and \emph{vacana}, \textquotedblleft{}speech, a word\textquotedblright{} (CDIAL 11199).
\end{cousinweb}

\begin{grammarlens}[title={What you've built}]
One more word for saying what happens if.
\end{grammarlens}

\subsection*{Your turn}
\begin{itemize}
\raggedright
  \item \emph{Say it:} \emph{pramāṇavacana}
  \item \emph{Say it:} \emph{pramāṇavacana}, once more
  \item \emph{Say it:} \emph{hosa mantrigaḷu pramāṇavacana tegedukoṇḍaru.} — The new ministers took the oath of office.
  \item \emph{Recall:} say the Kannada for a reward, then the Kannada for to ask formally, then say \emph{pramāṇavacana} again
\end{itemize}

\begin{tcolorbox}[breakable,colback=teal!4,colframe=teal!35!black,title={Before you move on}]
What does \kn{ಪ್ರಮಾಣವಚನ} mean? (\textquotedblleft{}An oath of office\textquotedblright{}.) Say it once more.
\end{tcolorbox}
